\chapter{Программная реализация} \label{ch5}

В данной главе описана программная реализация алгоритма GA-FP+Div. 
В параграфе~\ref{ch5:sec1} описана структура проекта. 
В параграфе~\ref{ch5:sec2} описаны ключевые классы. 
В параграфе~\ref{ch5:sec3} описаны форматы входных данных. 
В параграфе~\ref{ch5:sec4} приведена инструкция по установке и запуску.

\section{Структура проекта} \label{ch5:sec1}

Программная реализация выполнена на языке Java с использованием системы сборки Maven. 
Проект организован в следующие пакеты:

\begin{itemize}
    \item \texttt{com.tcp.algorithm} --- алгоритмы приоритизации: 
    \texttt{AdditionalPrioritizer}, \texttt{GeneticPrioritizer}, \texttt{RandomPrioritizer}, интерфейс \texttt{Prioritizer};
    \item \texttt{com.tcp.fitness} --- фитнесс-функция: \texttt{GaFitnessFunction};
    \item \texttt{com.tcp.data} --- загрузка данных: \texttt{DataLoader};
    \item \texttt{com.tcp.metric} --- метрики оценки: \texttt{ApfdCalculator}, \texttt{ApfdCostCalculator};
    \item \texttt{com.tcp.model} --- модель данных: \texttt{TestCase};
    \item \texttt{com.tcp} --- точка входа: \texttt{Main}.
\end{itemize}

\section{Описание ключевых классов} \label{ch5:sec2}

\textbf{Интерфейс \texttt{Prioritizer}} определяет единый программный интерфейс для всех алгоритмов 
приоритизации: метод \texttt{prioritize(List<TestCase>)} принимает неупорядоченный пул тестов и 
возвращает упорядоченный список.

\textbf{Класс \texttt{TestCase}} представляет тестовый случай и содержит четыре поля: идентификатор теста, 
множество покрываемых методов, вероятность дефектов и время выполнения.

\textbf{Класс \texttt{GaFitnessFunction}} реализует многокомпонентную фитнесс-функцию. 
Метод \texttt{calculateFitness(List<TestCase>)} вычисляет взвешенную сумму трёх компонент: покрытия, 
вероятности дефектов и структурного разнообразия. Веса компонент задаются в конструкторе.

\textbf{Класс \texttt{GeneticPrioritizer}} реализует генетический алгоритм. Конструктор принимает размер 
популяции, число поколений, вероятности кроссовера и мутации, а также экземпляр \texttt{GaFitnessFunction}.
Метод \texttt{prioritize} выполняет эволюционный цикл и возвращает лучшую найденную перестановку.

\textbf{Класс \texttt{DataLoader}} обеспечивает загрузку данных трёх форматов: BigFaultMatrix, Defects4J 
и синтетических данных. Все форматы приводятся к единому внутреннему представлению \texttt{FaultMatrixData},
содержащему список тестов, карту убитых мутантов и общее число мутантов.

\textbf{Класс \texttt{ApfdCalculator}} реализует расчёт метрики APFD по формуле~(\ref{eq:apfd}). 

\textbf{Класс \texttt{ApfdCostCalculator}} реализует расчёт метрики APFDc с учётом времени выполнения 
тестов.

\section{Форматы входных данных} \label{ch5:sec3}

Программа поддерживает три формата входных данных.

\textbf{BigFaultMatrix} --- бинарная матрица тест~$\times$~дефект в текстовом формате. 
Каждая строка соответствует одному тесту: первый элемент --- идентификатор теста, остальные --- значения 0 
или 1, где 1 означает что тест обнаруживает соответствующий дефект.

\begin{lstlisting}[caption={Фрагмент BigFaultMatrix}]
t1,0,1,0,1,0,0,1,...
t2,1,0,1,0,1,0,0,...
\end{lstlisting}

\textbf{Defects4J} --- данные реальных Java-проектов. Для каждого проекта используются три файла:
\begin{itemize}
    \item \texttt{coverage/matrix.csv} --- матрица покрытия методов (строки: тесты, столбцы: методы);
    \item \texttt{mutants/kill.csv}, \texttt{testMap.csv}, \texttt{covMap.csv} --- информация об убитых мутантах;
    \item \texttt{cia/impact\_probabilities.csv} --- вероятности дефектов методов на основе CIA-анализа.
\end{itemize}

\textbf{Синтетические данные} генерируются программно методом 

\texttt{DataLoader.generateSynthetic(numTests, numFaults, density)}, 
принимающим три параметра: количество тестов, количество дефектов и плотность покрытия. 
Для обеспечения нетривиального сравнения алгоритмов тесты делятся на три группы:

\begin{itemize}
    \item Тесты с широким покрытием и низким риском (30\% тестов) --- тесты с широким покрытием обычных дефектов, но низкой 
    вероятностью дефектов ($r_i \in [0{,}05;\ 0{,}10]$). Жадный алгоритм склонен выбирать их первыми 
    из-за большого покрытия, однако они не покрывают наиболее рискованные методы.

    \item Тесты с узким покрытием и высоким риском (20\% тестов) --- тесты с узким покрытием, но высокой вероятностью дефектов 
    ($r_i \in [0{,}60;\ 1{,}00]$). Покрывают первые 40\% дефектов, считающихся рискованными. 
    Алгоритм GA-FP+Div находит их раньше за счёт компоненты вероятности дефектов.

    \item Обычные тесты (50\% тестов) --- тесты со случайным покрытием с заданной плотностью 
    \texttt{density}. Вероятность дефектов пропорциональна доле покрытых дефектов.
\end{itemize}

Такая структура позволяет наглядно продемонстрировать различие в стратегиях алгоритмов: 
Additional выбирает тесты с максимальным покрытием методов, тогда как GA-FP+Div ставит в приоритет 
тесты с высокой вероятностью дефектов.

Результаты экспериментов сохраняются в двух форматах: 
\texttt{results/results.csv} для анализа в табличных редакторах и 
\texttt{results/results.json} для программной обработки.

\section{Инструкция по установке и запуску} \label{ch5:sec4}

\textbf{Требования:} Java 11 или выше, Apache Maven 3.6 или выше.

\textbf{Сборка проекта:}
\begin{lstlisting}
mvn clean package
\end{lstlisting}

\textbf{Запуск:}
\begin{lstlisting}
java -jar target/tcp-1.0.jar
\end{lstlisting}

После запуска отображается интерактивное меню выбора режима работы:
\begin{lstlisting}
=== TCP Benchmark ===
1. BigFaultMatrix (4 вариации)
2. Defects4J
3. Синтетические данные
4. Все Defects4J проекты
0. Выход
\end{lstlisting}

Для запуска на данных Defects4J необходимо указать путь к папке проекта, содержащей подпапки 
\texttt{coverage}, \texttt{mutants} и \texttt{cia}. Результаты автоматически сохраняются в папку 
\texttt{results/} при выходе из программы.

\textbf{Эталонный прогон.} Для проверки корректности работы программы предусмотрен запуск на наборе 
данных BigFaultMatrix (пункт меню 1). Ожидаемые результаты для оригинальной матрицы:

\begin{table}[h]
\caption{Эталонные результаты на BigFaultMatrix}
\label{tab:reference}
\begin{tabularx}{\textwidth}{|X|c|c|}
\hline
Алгоритм & APFD & APFDc \\
\hline
Random      & 0{,}9230 & 0{,}6861 \\
Additional  & 0{,}9987 & 0{,}7617 \\
GA-Base     & 0{,}9168 & 0{,}6799 \\
GA-FP+Div   & 0{,}9590 & 0{,}7221 \\
\hline
\end{tabularx}
\end{table}